% Chapter 8 of Dissolving Problems: Fourier analysis on the Boolean cube.
\newcommand{\fourierStab}{\operatorname{Stab}}
\newcommand{\fourierInf}{\operatorname{Inf}}
\newcommand{\fourierTI}{\mathbf{I}}
\newcommand{\fourierMaj}{\operatorname{Maj}}
\newcommand{\fourierW}[1]{\mathrm{W}^{#1}}
\newcommand{\fourierLin}{\mathrm{Lin}}
\newcommand{\fourierOnes}{\mathbf{1}}

\casechapter{Change of Basis}{Fourier analysis on the Boolean cube}{Diagonalize the symmetry}{ch:fourier}
\begin{refsection}[chapters/08-fourier.bib]

\begin{chaptersummary}
Condorcet's paradox, Arrow's impossibility theorem, and the linearity test of Blum, Luby, and Rubinfeld were settled with unrelated tools: central-limit asymptotics, voting profiles built by hand, and a self-correcting vote. Each asks for the probability of an event about one Boolean function at a few correlated random points whose joint law is invariant under flips of coordinates. In the characters of the cube, the basis that diagonalizes this symmetry, each probability becomes a weighted sum of powers of Fourier coefficients. Arrow's theorem and Guilbaud's $8.77\%$ become two readings of one formula, the linearity test gets a three-line analysis valid at every error rate, and a question about the power of single voters that counting could not settle yields to hypercontractivity.
\end{chaptersummary}

\section{Before: votes and tests, counted by hand}\label{sec:fourier:before}

The problems of this chapter come from social choice, which asks how a society should combine the rankings of its members, and from theoretical computer science, which asks how a few random queries can certify that a program is linear. We present three problems in their native settings, and then a fourth that neither toolkit could reach.

\subsection{Condorcet's paradox, and how often it occurs}\label{sec:fourier:condorcet}

Three voters rank three candidates $a$, $b$, $c$:
\[
  a\succ b\succ c,\qquad b\succ c\succ a,\qquad c\succ a\succ b .
\]
In the contest between $a$ and $b$, the first and third voters prefer $a$, so $a$ wins two votes to one. In the same way $b$ beats $c$ and $c$ beats $a$ (\cref{fig:fourier:cycle}): majority rule applied to each pair produces a cycle, although every voter is consistent. \Textcite{Condorcet1785} described this phenomenon, now called \term{Condorcet's paradox}\index{Condorcet's paradox}. A candidate who beats each of the others in pairwise majority votes is a \term{Condorcet winner}; with three candidates and an odd number of voters there is one exactly when the three pairwise outcomes do not form a cycle.

\begin{figure}[t]
\centering
\begin{tikzpicture}[>={Stealth[length=2.6mm]}]
  \tikzset{cand/.style={circle,draw=sea,fill=seafoam,line width=0.6pt,minimum size=7.5mm,inner sep=0pt}}
  \node[cand] (a) at (90:1.3) {$a$};
  \node[cand] (b) at (210:1.3) {$b$};
  \node[cand] (c) at (330:1.3) {$c$};
  \draw[->,sea,line width=0.8pt] (a) -- node[left=1pt,font=\footnotesize,text=ink] {$2$--$1$} (b);
  \draw[->,sea,line width=0.8pt] (b) -- node[below=1pt,font=\footnotesize,text=ink] {$2$--$1$} (c);
  \draw[->,sea,line width=0.8pt] (c) -- node[right=1pt,font=\footnotesize,text=ink] {$2$--$1$} (a);
  \node[anchor=west,font=\small] at (3.0,0.5) {voter 1:\enspace $a\succ b\succ c$};
  \node[anchor=west,font=\small] at (3.0,0.0) {voter 2:\enspace $b\succ c\succ a$};
  \node[anchor=west,font=\small] at (3.0,-0.5) {voter 3:\enspace $c\succ a\succ b$};
\end{tikzpicture}
\caption{Condorcet's paradox. Each arrow points from the winner to the loser of a pairwise majority vote among the three voters on the right.}
\label{fig:fourier:cycle}
\end{figure}

How often does this happen? Under the model now called \term{impartial culture}\index{impartial culture}, each of $n$ voters picks one of the six rankings uniformly at random, independently of the others. With three voters, exactly $12$ of the $6^3=216$ profiles produce a cycle, namely the three rankings above in some order or the three reverse rankings in some order, so the probability is $1/18$. For majority rule with an odd number $n$ of voters, \textcite{Guilbaud1952} stated the limit
\begin{equation}\label{eq:fourier:guilbaud}
  \lim_{n\to\infty}\Pr[\text{cycle}]=\frac14-\frac{3}{2\pi}\arcsin\frac13\approx0.0877,
\end{equation}
with little explanation of its derivation.\sidenote{Guilbaud's long paper appeared in the same 1952 issue of \emph{\'Economie appliqu\'ee} as a French version of a lecture by Arrow. A complete English translation appeared only in 2008.}

With tools available in 1952 one argues as follows. Encode the ranking of voter $i$ by three signs: $x_i=+1$ if the voter ranks $a$ above $b$ and $x_i=-1$ otherwise, and similarly $y_i$ for $b$ against $c$ and $z_i$ for $c$ against $a$. Each sign is uniform, and a count over the six rankings gives $\E[x_iy_i]=\E[y_iz_i]=\E[z_ix_i]=-\frac13$. By the central limit theorem the margins $n^{-1/2}(\sum_ix_i,\sum_iy_i,\sum_iz_i)$ converge in law to a centered Gaussian vector $(X,Y,Z)$ with unit variances and correlations $-\frac13$, which puts no mass on the coordinate planes, and a cycle is the event that the three margins have one sign. Sheppard's quadrant probability $\Pr[X>0,Y>0]=\frac14+\frac1{2\pi}\arcsin\rho$ for jointly Gaussian standard variables with correlation $\rho$ \parencite{Sheppard1899}, with inclusion and exclusion over the three pairs and the symmetry $\Pr[X,Y,Z<0]=\Pr[X,Y,Z>0]$, gives $\Pr[X,Y,Z>0]=\frac18-\frac{3}{4\pi}\arcsin\frac13$, and doubling for the two orientations gives \eqref{eq:fourier:guilbaud}. The computation is tied to majority and to $n\to\infty$: another rule needs another limit theorem, five voters need $6^5$ profiles counted, and nothing hints at a connection with the next problem.

\subsection{Arrow's theorem: never?}\label{sec:fourier:arrow}

\Textcite{Arrow1951} asked whether some rule for combining individual rankings into a social ranking avoids Condorcet's paradox while respecting a few mild conditions. We state his answer for three candidates and strict rankings, the case that carries the difficulty.\sidenote{Arrow's theorem holds for every finite number of candidates at least three, and Arrow allowed ties. The first edition of his book stated it under conditions too weak for it to hold, as \textcite{Blau1957} showed by a counterexample; the second edition corrected this.}

The central condition, \term{independence of irrelevant alternatives}\index{independence of irrelevant alternatives}, asks that society's verdict on two candidates depend only on the voters' preferences between those two. With three candidates, such a rule is given by three functions $f,g,h\colon\{-1,1\}^n\to\{-1,1\}$: society ranks $a$ above $b$ when $f(x)=1$, $b$ above $c$ when $g(y)=1$, and $c$ above $a$ when $h(z)=1$, with $x$, $y$, $z$ the sign vectors above. A \term{profile} is a triple $(x,y,z)$ of sign vectors such that for each voter $i$ the signs $x_i$, $y_i$, $z_i$ are not all equal; these are exactly the triples that come from rankings. The rule $(f,g,h)$ is \term{rational} if it never produces a cycle, that is, if $f(x)$, $g(y)$, $h(z)$ are not all equal for each profile; \term{unanimous} if $f$, $g$, $h$ map the all-ones vector $\fourierOnes$ to $1$ and $-\fourierOnes$ to $-1$; and a \term{dictator} if $f=g=h=\chi_{\{i\}}$ for some voter $i$, with $\chi_{\{i\}}(x)=x_i$.

\begin{theorem}[Arrow's theorem for three candidates]\label{thm:fourier:arrow}\index{Arrow's theorem}
Every unanimous rational rule on three candidates is a dictator.
\end{theorem}

The combinatorial proofs share a shape: build a few profiles by hand, read off constraints, and assemble them into a dictator. The aside carries this out when one function $f$ decides all three contests, a case to which the general one reduces (Exercise~\ref{ex:fourier:neutral}).

\begin{aside}[A combinatorial proof, one rule for all three contests]
Let $f$ be rational. For $C\subseteq[n]$ let $x^C$ be the sign vector that is $+1$ on $C$ and $-1$ off $C$. Each step exhibits a profile; the reader checks that no voter's three signs are equal.
\emph{Oddness.} The profile $(x,-x,x)$ forces $f(-x)\ne f(x)$, so $f(-x)=-f(x)$. Replacing $f$ by $-f$ if necessary, which preserves rationality, assume $f(\fourierOnes)=1$.
\emph{Upward closure.} If $f(x^A)=1$ and $A\subseteq D$, then the profile $(x^D,-x^A,-\fourierOnes)$ has outcomes $(f(x^D),-1,-1)$, so $f(x^D)=1$.
\emph{Splitting.} If $f(x^C)=1$ and $C$ is the disjoint union of $A$ and $B$, then the profile $(x^C,-x^A,-x^B)$ has outcomes $(1,-f(x^A),-f(x^B))$, so $f(x^A)=1$ or $f(x^B)=1$.
\emph{The dictator.} The family $\cD=\{C\colon f(x^C)=1\}$ contains $[n]$ and, by oddness, exactly one set of each complementary pair. A smallest member $C$ of $\cD$ is nonempty, since $\emptyset\notin\cD$, and splitting forbids it to have two elements or more; so $C=\{i\}$. Upward closure puts every set containing $i$ into $\cD$, and oddness keeps every other set out. Hence $f(x)=x_i$ for every $x$; before the normalization, $f=\pm\chi_{\{i\}}$.
\end{aside}

The argument is short and clever, and it answers exactly the question asked: for each unanimous rule that is not a dictator it exhibits one bad profile, that is, a cycle with probability at least $6^{-n}$ under impartial culture. It cannot tell whether a rule far from every dictator must produce cycles \emph{often}, a question raised by \textcite{Kalai2002}, and it is silent about majority's $8.77\%$.

\subsection{Linearity testing: what does passing prove?}\label{sec:fourier:blr}

\Textcite{BlumLubyRubinfeld1993}, in work presented in 1990, asked how a program that claims to compute a linear map can be checked on a few random inputs. For $f\colon\F_2^n\to\F_2$ their test picks $x,y\in\F_2^n$ independently and uniformly and accepts if $f(x)+f(y)=f(x+y)$.\index{linearity testing} The linear functions $x\mapsto\sum_{i\in S}x_i$, for $S\subseteq[n]$, always pass. Write $\delta(f)$ for the probability that the test rejects $f$, and $\dist(f,\fourierLin)$ for the smallest fraction of inputs on which $f$ differs from a linear function. How does $\delta(f)$ control $\dist(f,\fourierLin)$? The original answer is a \term{self-correction}\index{self-correction} argument, in which the values of $f$ vote for the linear function that $f$ is trying to be.

\begin{aside}[The self-correction argument, in a streamlined form]
Let $\delta=\delta(f)$, and let $g(x)$ be the more frequent value of $f(x+y)+f(y)$ as $y$ ranges over $\F_2^n$, with ties broken arbitrarily.
\emph{The vote is close to $f$.} If $f(x)\ne g(x)$, then the test with first query $x$ rejects for at least half of the $y$. Averaging over $x$ gives $\delta\ge\frac12\Pr_x[f(x)\ne g(x)]$.
\emph{Every vote is a landslide.} Fix $x$ and draw $y$, $z$ independently. The pairs $(x+y,z)$ and $(y,x+z)$ are uniform, so with probability at least $1-2\delta$ both tests pass, and then $f(x+y)+f(z)=f(x+y+z)=f(y)+f(x+z)$: the votes of $y$ and of $z$ agree. If a fraction $p\ge\frac12$ of the votes go to the majority value, then two independent votes agree with probability $p^2+(1-p)^2\le p$, so $p\ge1-2\delta$.
\emph{The vote is linear if $\delta<1/6$.} Fix $x$ and $x'$ and draw $y$. By the landslide property each of the identities $g(x)=f(x+y)+f(y)$, $g(x')=f(x+x'+y)+f(x+y)$, and $g(x+x')=f(x+x'+y)+f(y)$ fails with probability at most $2\delta$, so for some $y$ all three hold; adding the first two gives $g(x)+g(x')=g(x+x')$.
\emph{Conclusion.} If $\delta(f)<1/6$, then $\dist(f,\fourierLin)\le2\delta(f)$.
\end{aside}

For maps between arbitrary finite groups, Blum, Luby, and Rubinfeld proved that the rejection probability is at least $\frac29$ of the distance, crediting Coppersmith with part of the analysis \parencite[as recounted by][]{BellareEtAl1996}, and \textcite{BellareEtAl1993} proved the bound $3d-6d^2$, much better for small $d=\dist(f,\fourierLin)$. These constants mattered, because each improvement gave a better hardness-of-approximation factor for satisfiability problems.\sidenote{Probabilistically checkable proofs use the test to check that a proof is a codeword of the Hadamard code, whose codewords are the linear functions \parencite{BellareEtAl1996}.} But the decisive step of the method needs $\delta<1/6$, and for a function that passes with probability $0.6$ these bounds say nothing at all, not even that $f$ is correlated with some linear function.

\subsection{A fourth question: how much power must one voter have?}\label{sec:fourier:power}

Let $f\colon\{-1,1\}^n\to\{-1,1\}$ decide an election between two candidates. Voter $i$ is \term{pivotal} at $x$ if flipping $x_i$ alone changes $f(x)$, and the \term{influence}\index{influence} $\fourierInf_i[f]$ is the probability that voter $i$ is pivotal when $x$ is uniform. A dictator gives one voter influence $1$; majority gives each voter the probability $\binom{n-1}{(n-1)/2}2^{1-n}\approx\sqrt{2/(\pi n)}$ that the other votes split evenly. Studying collective coin flipping, \textcite{BenOrLinial1985} asked how small the largest influence can be when $f$ is \term{balanced}, taking each value with probability $\frac12$. The edge-isoperimetric inequality of the cube makes the average influence at least $1/n$. Ben-Or and Linial constructed the \term{tribes} function\index{tribes function}, in which every voter has influence of order $(\log n)/n$, and conjectured that this is optimal up to a constant factor. When \textcite{KahnKalaiLinial1988} proved the conjecture, they remarked that no purely combinatorial argument had improved on $1/n$.

Four questions, four toolkits: limit theorems, hand-built profiles, a self-correcting vote, and edge counting.

\section{The move: diagonalize the symmetry}\label{sec:fourier:move}

Look again at the four questions. Each evaluates one function on the cube at a few random points, namely $x$, $y$, $z$ for Condorcet and Arrow, $x$, $y$, $x+y$ for the test, and $x$ and $x$ with a coordinate flipped for influence, and asks for the expectation of a polynomial in the values. In each, the joint law of the points is invariant under flipping coordinates: reversing one voter's ballot flips that voter's coordinate in $x$, $y$, and $z$ at once and maps impartial culture to itself, and flipping a coordinate in the first and third queries of the test, or in both points of the influence question, preserves their law. The move is to work in the basis in which these symmetries act most simply.

Identify $\F_2^n$ with $\{-1,1\}^n$ coordinatewise through $b\mapsto(-1)^b$; addition in $\F_2^n$ becomes the coordinatewise product, written $x\odot y$. Let $x$ be uniform on the cube, write $\E$ for expectation with respect to it, and give the $2^n$-dimensional space of functions $f\colon\{-1,1\}^n\to\R$ the inner product $\ip{f,g}=\E[f(x)g(x)]$ and the norms $\norm{f}_p=\E[\abs{f(x)}^p]^{1/p}$. A function is \term{Boolean} if it takes only the values $\pm1$.

\begin{definition}[Characters and Fourier coefficients]\label[definition]{def:fourier:characters}
For $S\subseteq[n]$, the \term{character}\index{character!of the cube} $\chi_S$ is the parity $\chi_S(x)=\prod_{i\in S}x_i$, with $\chi_\emptyset=1$. The \term{Fourier coefficient} of $f$ at $S$ is $\hat f(S)=\ip{f,\chi_S}$, and the \term{weight of $f$ at level $k$} is $\fourierW{k}[f]=\sum_{\abs S=k}\hat f(S)^2$.
\end{definition}

In the $\pm1$ encoding the linear functions become exactly the characters, the maps that turn the group operation into multiplication: $\chi_S(x\odot y)=\chi_S(x)\chi_S(y)$.\sidenote{The functions $\chi_S$ go back to \textcite{Walsh1923}, who introduced the corresponding orthonormal system on the unit interval; through binary expansions his functions are the characters of the cube.}

\begin{theorem}[Fourier--Walsh expansion]\label{thm:fourier:walsh}\index{Fourier--Walsh expansion}\index{Parseval's identity}
The characters form an orthonormal basis of the functions on $\{-1,1\}^n$. Consequently every $f$ equals $\sum_S\hat f(S)\chi_S$; every $f$ and $g$ satisfy $\ip{f,g}=\sum_S\hat f(S)\hat g(S)$ (Plancherel's identity); and $\E[f(x)^2]=\sum_S\hat f(S)^2$ (Parseval's identity), which equals $1$ for Boolean $f$.
\end{theorem}

\begin{proof}
Because $x_i^2=1$, the product $\chi_S\chi_T$ equals $\chi_{S\triangle T}$, so by independence of the coordinates $\ip{\chi_S,\chi_T}=\prod_{i\in S\triangle T}\E[x_i]$. This product is $1$ if $S=T$, being empty, and $0$ otherwise. The $2^n$ characters are therefore orthonormal, hence a basis of the $2^n$-dimensional space, and the three identities are the usual expansions in an orthonormal basis.
\end{proof}

Why characters? Because they are the eigenvectors of the symmetry. For $a\in\{-1,1\}^n$ let $\tau_a$ be the translation $(\tau_af)(x)=f(a\odot x)$.

\begin{lemma}[Invariant operators are diagonal]\label[lemma]{lem:fourier:diag}
If a linear operator $K$ on functions on $\{-1,1\}^n$ commutes with every translation $\tau_a$, then every character is an eigenvector of $K$.
\end{lemma}

\begin{proof}
Each character is an eigenvector of each translation, because $\chi_S(a\odot x)=\chi_S(a)\chi_S(x)$. Expand $K\chi_S=\sum_Tc_T\chi_T$. Commutation gives $\tau_aK\chi_S=K\tau_a\chi_S=\chi_S(a)K\chi_S$, and comparing coefficients yields $c_T\chi_T(a)=c_T\chi_S(a)$ for every $a$ and $T$. For $T\ne S$, choose $a$ equal to $-1$ at one coordinate of $S\triangle T$ and $+1$ elsewhere; then $\chi_T(a)=-\chi_S(a)$, which forces $c_T=0$. Hence $K\chi_S=c_S\chi_S$.
\end{proof}

\begin{definition}[Correlated pairs, noise, stability]\label[definition]{def:fourier:noise}
Let $-1\le\rho\le1$. A pair $(x,y)$ of random points of the cube is \term{$\rho$-correlated} if the pairs $(x_1,y_1),\dots,(x_n,y_n)$ are independent and each satisfies $\E[x_i]=\E[y_i]=0$ and $\E[x_iy_i]=\rho$; equivalently, $x$ is uniform and each $y_i$ equals $x_i$ with probability $\frac{1+\rho}2$, independently. The \term{noise operator}\index{noise operator} is $T_\rho f(x)=\E[f(y)\mid x]$, and the \term{noise stability}\index{noise stability} of $f$ is $\fourierStab_\rho[f]=\E[f(x)f(y)]$.
\end{definition}

\begin{lemma}[The noise operator is diagonal]\label[lemma]{lem:fourier:noise}
For each $S$, $T_\rho\chi_S=\rho^{\abs S}\chi_S$. Consequently, for a $\rho$-correlated pair $(x,y)$ and all functions $f$ and $g$,
\begin{equation}\label{eq:fourier:noisecorr}
  \E[f(x)g(y)]=\sum_{S\subseteq[n]}\rho^{\abs S}\hat f(S)\hat g(S);
\end{equation}
in particular $\fourierStab_\rho[f]=\sum_{k=0}^n\rho^k\fourierW{k}[f]$.
\end{lemma}

\begin{proof}
Multiplying both coordinates of a pair $(x_i,y_i)$ by the same sign preserves the means and the correlation, so the law of $(x,y)$ is invariant under $(x,y)\mapsto(a\odot x,a\odot y)$. Hence $T_\rho$ commutes with translations, and \cref{lem:fourier:diag} predicts that the characters are its eigenvectors. One computation gives the eigenvalues: given $x$, the coordinates $y_i$ are independent with $\E[y_i\mid x]=\rho x_i$, so $T_\rho\chi_S(x)=\prod_{i\in S}\rho x_i=\rho^{\abs S}\chi_S(x)$. For the identity, write $\E[f(x)g(y)]=\ip{f,T_\rho g}$, expand $T_\rho g=\sum_S\rho^{\abs S}\hat g(S)\chi_S$, and apply Plancherel's identity.
\end{proof}

These two facts are the whole dictionary: a probability about $f$ at correlated points becomes a sum over $S$ of an eigenvalue times a product of Fourier coefficients, and Parseval's identity supplies the one constraint that Booleanness imposes. \Cref{tab:fourier:dictionary} collects the entries we use; the agreement row follows from $\E[f\chi_S]=\Pr[f=\chi_S]-\Pr[f\ne\chi_S]$, and the last two rows are proved below.

\begin{table}[b]
\centering\small
\caption{The dictionary of this chapter. Here $(x,y)$ is a $\rho$-correlated pair; in the linearity-test row $\rho=0$, so $x$ and $y$ are independent. The last three rows assume that $f$ is Boolean.}
\label{tab:fourier:dictionary}
\begin{tabular}{@{}lll@{}}
\toprule
Quantity & In values & In Fourier coefficients\\
\midrule
mean & $\E f(x)$ & $\hat f(\emptyset)$\\
variance & $\Var f(x)$ & $\sum_{S\ne\emptyset}\hat f(S)^2$\\
correlation under noise & $\E[f(x)g(y)]$ & $\sum_S\rho^{\abs S}\hat f(S)\hat g(S)$\\
agreement with $\chi_S$ & $\Pr[f(x)=\chi_S(x)]$ & $\frac12\bigl(1+\hat f(S)\bigr)$\\
linearity test & $\Pr[f(x)f(y)=f(x\odot y)]$ & $\frac12+\frac12\sum_S\hat f(S)^3$\\
influence of voter $i$ & $\Pr[f(x)\ne f(x^{\oplus i})]$ & $\sum_{S\ni i}\hat f(S)^2$\\
\bottomrule
\end{tabular}
\end{table}

The characters are parities, the Boolean functions a combinatorialist would rank as the least interesting: each is orthogonal to every function that ignores one of its variables. Grothendieck's advice to stoop low applies literally. Parities keep their shape under the symmetry, since translating a parity only multiplies it by a sign, so every operator built from a symmetric experiment acts on them by a number. What one had to notice in the four problems was neither a clever profile nor a clever vote, but a symmetry that the random experiment has for free, together with the decision to stop looking at the values of $f$ and look at its spectrum.\sidenote{The same diagonalization, applied to the operator of a graphon rather than to a group, gives the spectral characterization of quasirandom graphs in \cref{ch:graphon}.}

\section{After: tests and elections on the spectrum}\label{sec:fourier:after}

\subsection{Linearity testing in three lines}\label{sec:fourier:blrafter}

In $\pm1$ notation the test of \cref{sec:fourier:blr} picks independent uniform points $x$ and $y$ and accepts if $f(x)f(y)=f(x\odot y)$; the linear functions are the characters, and $\dist(f,\fourierLin)=\min_S\Pr[f\ne\chi_S]$.

\begin{theorem}[\textcite{BellareEtAl1996}]\label{thm:fourier:blr}
Every Boolean function $f$ on $\{-1,1\}^n$ passes the linearity test with probability
\[
  \frac12+\frac12\sum_{S\subseteq[n]}\hat f(S)^3 .
\]
Consequently, if $f$ passes with probability $\frac12+\eta$, then $f$ agrees with some character on a fraction at least $\frac12+\eta$ of the points. In particular the test rejects $f$ with probability at least $\dist(f,\fourierLin)$.
\end{theorem}

\begin{proof}
Since $f(x)f(y)f(x\odot y)=\pm1$, the indicator of acceptance is $\frac12\bigl(1+f(x)f(y)f(x\odot y)\bigr)$. Expand the three factors in characters, using $\chi_U(x\odot y)=\chi_U(x)\chi_U(y)$ for the third:
\[
  \E[f(x)f(y)f(x\odot y)]=\sum_{S,T,U}\hat f(S)\hat f(T)\hat f(U)\,\E[\chi_S(x)\chi_U(x)]\,\E[\chi_T(y)\chi_U(y)].
\]
Orthonormality kills every term except $S=T=U$, which leaves $\sum_U\hat f(U)^3$. For the consequence, bound one factor of each cube by the largest coefficient and apply Parseval's identity: $\sum_U\hat f(U)^3\le\max_S\hat f(S)\sum_U\hat f(U)^2=\max_S\hat f(S)$. Hence $\hat f(S)\ge2\eta$ for some $S$, and the agreement of $f$ with $\chi_S$ is $\frac12\bigl(1+\hat f(S)\bigr)\ge\frac12+\eta$. Writing $\frac12+\eta=1-\delta(f)$ turns this into $\Pr[f\ne\chi_S]\le\delta(f)$.
\end{proof}

Set the two proofs side by side. The self-correction argument needs $\delta<1/6$ and loses a factor of $2$; the spectral argument needs no hypothesis and gives the constant $1$. It is shorter for a reason: the old argument must construct the linear function and check three properties of it, while the new one constructs nothing, because acceptance is a cubic form that the symmetry diagonalizes, and such a form with weights summing to $1$ is largest when all the weight sits on one coordinate. The regime the old argument could not touch is the important one. If $f$ passes with probability $\frac12+\eta$, then by Parseval's identity at most $1/(4\eta^2)$ characters have a coefficient as large as $2\eta$, so the test pins $f$ down to a short list of linear functions. The optimal inapproximability results of \textcite{Hastad2001}, among them the hardness of beating a random assignment for linear equations over $\F_2$ with three variables each, rest on Fourier analyses of tests of this kind with noise added.

\subsection{Condorcet and Arrow in one formula}\label{sec:fourier:kalai}

Return to three candidates and $n$ voters under impartial culture, with the sign vectors $x$, $y$, $z$ of \cref{sec:fourier:condorcet}.

\begin{lemma}[Impartial culture in coordinates]\label[lemma]{lem:fourier:culture}
Under impartial culture, each of the pairs $(x,y)$, $(y,z)$, and $(z,x)$ is $(-\frac13)$-correlated.
\end{lemma}

\begin{proof}
A ranking never yields three equal signs, and its signs determine all pairwise comparisons, hence the ranking. So the map from rankings to sign triples is a bijection onto the six triples in $\{-1,1\}^3$ that are not constant, and each voter's triple is uniform on these six, independently across voters. The six triples are closed under negation, so each sign has mean $0$; and for each two of the three positions, exactly two of the six triples agree there, so each product of two signs has mean $(2-4)/6=-\frac13$.
\end{proof}

\begin{theorem}[Kalai's formula; \textcite{Kalai2002}]\label{thm:fourier:kalai}\index{Kalai's formula}
Under impartial culture, a rule that uses the Boolean function $f$ in all three contests produces a cycle with probability
\begin{equation}\label{eq:fourier:kalai}
  \Pr[\text{cycle}]=\frac14+\frac34\fourierStab_{-1/3}[f]=\frac14+\frac34\sum_{k=0}^n\Bigl(-\frac13\Bigr)^k\fourierW{k}[f].
\end{equation}
\end{theorem}

\begin{proof}
For signs $w_1$, $w_2$, $w_3$ the number $\frac14(1+w_1w_2+w_2w_3+w_3w_1)$ is $1$ if the three are equal and $0$ otherwise, since otherwise exactly one of the three products is $+1$. A cycle is the event $f(x)=f(y)=f(z)$, so taking expectations,
\[
  \Pr[\text{cycle}]=\frac14+\frac14\bigl(\E[f(x)f(y)]+\E[f(y)f(z)]+\E[f(z)f(x)]\bigr).
\]
By \cref{lem:fourier:culture} each expectation is $\fourierStab_{-1/3}[f]$, which \eqref{eq:fourier:noisecorr} evaluates.
\end{proof}

\begin{example}[Three voters]\label[example]{exm:fourier:maj3}
Checking the four kinds of inputs shows $\fourierMaj_3(x)=\frac12(x_1+x_2+x_3)-\frac12x_1x_2x_3$. So $\fourierW1=\frac34$ and $\fourierW3=\frac14$, and \eqref{eq:fourier:kalai} gives $\frac14+\frac34\bigl(-\frac14-\frac1{108}\bigr)=\frac1{18}$, as the count of $12$ profiles out of $216$ did. For five voters the formula gives $\frac5{72}$ (Exercise~\ref{ex:fourier:maj5}), with no profile counted.
\end{example}

Arrow's theorem asks when \eqref{eq:fourier:kalai} vanishes. The weight must then sit on level $1$, and the next lemma says what that forces.

\begin{lemma}[Weight on level one]\label[lemma]{lem:fourier:level1}
If a Boolean function $f$ has $\fourierW1[f]=1$, then $f=\chi_{\{i\}}$ or $f=-\chi_{\{i\}}$ for some $i$.
\end{lemma}

\begin{proof}
By Parseval's identity $f$ has no weight off level $1$, so $f(x)=\sum_ja_jx_j$ with $a_j=\hat f(\{j\})$ and $\sum_ja_j^2=1$. Flipping the $j$th coordinate of $x$ changes $f(x)$ by $-2a_jx_j$, and since $f$ takes only the values $\pm1$, this change lies in $\{-2,0,2\}$. So each $a_j$ is $-1$, $0$, or $1$, and as $\sum_ja_j^2=1$, exactly one of them is nonzero.
\end{proof}

\begin{theorem}[Arrow's theorem for one rule, with a number]\label{thm:fourier:arrowq}
Under impartial culture, a rule that uses the Boolean function $f$ in all three contests produces a cycle with probability at least $\frac29\bigl(1-\fourierW1[f]\bigr)$. In particular, if $f$ never produces a cycle, then $f=\pm\chi_{\{i\}}$ for some $i$. If moreover $f(\fourierOnes)=1$, then $f$ is the dictator $\chi_{\{i\}}$.
\end{theorem}

\begin{proof}
The plan is to bound the sum in \eqref{eq:fourier:kalai} level by level, with Parseval's identity as the only constraint. Even levels contribute nonnegative terms. Level $1$ contributes $-\frac13\fourierW1[f]$, and each odd level $k\ge3$ contributes at least $-\frac1{27}\fourierW{k}[f]$. The levels other than the first carry total weight $1-\fourierW1[f]$, so
\[
  \Pr[\text{cycle}]\ge\frac14-\frac34\cdot\frac13\,\fourierW1[f]-\frac34\cdot\frac1{27}\bigl(1-\fourierW1[f]\bigr)=\frac29\bigl(1-\fourierW1[f]\bigr).
\]
If $f$ never produces a cycle, then this probability is $0$, so $\fourierW1[f]=1$, and \cref{lem:fourier:level1} gives $f=\pm\chi_{\{i\}}$. The condition $f(\fourierOnes)=1$ excludes the minus sign.
\end{proof}

Compare this proof with the aside of \cref{sec:fourier:arrow}: the three hand-built profiles and the family $\cD$ give way to one identity, which averages over all profiles at once, and one inequality, which reads off the spectrum. The combinatorial proof is short too, and as \textcite{Keller2012} remarks, an analytic proof of Arrow's theorem is not the point in itself. The point is the number. The combinatorial proof certifies only a probability of at least $6^{-n}$; \cref{thm:fourier:arrowq} gives a probability bounded away from $0$ independently of $n$ unless $\fourierW1[f]$ is close to $1$, and for majority on three voters the bound $\frac29\cdot\frac14=\frac1{18}$ is attained.

\begin{remark}[Three different rules]\label[remark]{rem:fourier:threerules}
For a rule $(f,g,h)$ the same proof gives $\Pr[\text{cycle}]=\frac14+\frac14\sum_S(-\frac13)^{\abs S}\bigl(\hat f(S)\hat g(S)+\hat g(S)\hat h(S)+\hat h(S)\hat f(S)\bigr)$. The spectral bound alone cannot force a dictator here, because constant rules can be rational: $f=g=1$ and $h=-1$ always rank $a\succ b\succ c$. Unanimity excludes such rules, and a short combinatorial argument shows that unanimity and rationality force $f=g=h$ (Exercise~\ref{ex:fourier:neutral}); \cref{thm:fourier:arrowq} then proves \cref{thm:fourier:arrow}.
\end{remark}

The quantitative form opened a line of work. \Textcite{FriedgutKalaiNaor2002} proved that a Boolean function with $\fourierW1[f]\ge1-\eps$ is within distance $O(\eps)$ of a dictator or its negation, and \textcite{Kalai2002} deduced that a neutral rule, one that treats the candidates symmetrically, which produces cycles with probability $\delta$ is within distance $O(\delta)$ of a dictator or its negation. General rules are much harder: with $\eps$ the distance to the rules that never cycle, \textcite{Mossel2012} proved that cycles have probability at least $\exp(-C/\eps^{21})$, and \textcite{Keller2012} improved this to $C\eps^3$, tight up to logarithmic factors.

\subsection{Guilbaud's number}\label{sec:fourier:guilbaud}

For majority, \eqref{eq:fourier:kalai} reduces Guilbaud's problem to the noise stability of $\fourierMaj_n(x)=\sgn(x_1+\dots+x_n)$, for odd $n$, and a two-dimensional Gaussian fact.

\begin{lemma}[Sheppard's formula]\label[lemma]{lem:fourier:sheppard}\index{Sheppard's formula}
If $(X,Y)$ is a centered Gaussian vector with unit variances and correlation $\rho\in[-1,1]$, then $\Pr[\sgn X\ne\sgn Y]=\frac1\pi\arccos\rho$; hence $\E[\sgn X\sgn Y]=\frac2\pi\arcsin\rho$.
\end{lemma}

\begin{proof}
Let $u$ and $v$ be unit vectors in $\R^2$ at angle $\theta=\arccos\rho$, and let $G$ be a standard Gaussian vector in $\R^2$; then $(\ip{u,G},\ip{v,G})$ has the law of $(X,Y)$. The direction of $G$ is uniform on the circle, and $\ip{u,G}$ and $\ip{v,G}$ have opposite signs exactly when the line through $0$ orthogonal to $G$ separates $u$ from $v$, which happens for a set of directions of total angle $2\theta$ out of $2\pi$. The second statement follows from $\E[\sgn X\sgn Y]=1-2\Pr[\sgn X\ne\sgn Y]$ and $\arccos\rho=\frac\pi2-\arcsin\rho$.
\end{proof}

\begin{theorem}[Stability of majority]\label{thm:fourier:majstab}
As $n\to\infty$ through odd values, $\fourierStab_\rho[\fourierMaj_n]\to\frac2\pi\arcsin\rho$ for each $\rho\in[-1,1]$.
\end{theorem}

\begin{proof}
Let $(x,y)$ be $\rho$-correlated. The vectors $(x_i,y_i)\in\R^2$ are independent and identically distributed with mean $0$ and covariance matrix $\bigl(\begin{smallmatrix}1&\rho\\\rho&1\end{smallmatrix}\bigr)$, so by the multivariate central limit theorem $n^{-1/2}\sum_i(x_i,y_i)$ converges in law to a pair $(X,Y)$ as in \cref{lem:fourier:sheppard}. The function $(s,t)\mapsto\sgn s\sgn t$ is bounded and continuous off the coordinate axes, which have probability $0$ for $(X,Y)$, so expectations converge: $\fourierStab_\rho[\fourierMaj_n]=\E\bigl[\sgn\bigl(\sum_ix_i\bigr)\sgn\bigl(\sum_iy_i\bigr)\bigr]\to\E[\sgn X\sgn Y]$.
\end{proof}

\begin{corollary}[Guilbaud's number]\label[corollary]{cor:fourier:guilbaud}
Under impartial culture with majority in every contest, $\Pr[\text{cycle}]\to\frac14-\frac3{2\pi}\arcsin\frac13\approx0.0877$ as the odd number of voters tends to infinity.
\end{corollary}

\begin{proof}
Put $\rho=-\frac13$ in \cref{thm:fourier:majstab}, substitute in \eqref{eq:fourier:kalai}, and use $\arcsin(-\frac13)=-\arcsin\frac13$.
\end{proof}

\begin{figure}[t]
\centering
\begin{tikzpicture}[x=0.152cm,y=0.3cm]
  \draw[muted,line width=0.4pt] (0,0) -- (63,0);
  \draw[muted,line width=0.4pt] (0,0) -- (0,10.3);
  \foreach \yy in {0,2,4,6,8,10} {\draw[muted,line width=0.4pt] (0,\yy) -- (-0.8,\yy) node[left,font=\scriptsize,text=ink] {$\yy\%$};}
  \foreach \xx in {1,11,21,31,41,51,61} {\draw[muted,line width=0.4pt] (\xx,0) -- (\xx,-0.25) node[below,font=\scriptsize,text=ink] {$\xx$};}
  \draw[sea,dashed,line width=0.6pt] (0,8.774) -- (63,8.774);
  \node[anchor=south east,font=\scriptsize,text=seadeep] at (63,8.9) {Guilbaud's limit $\approx 8.77\%$};
  \draw[seadeep,line width=0.6pt] plot coordinates {(1,0.000) (3,5.556) (5,6.944) (7,7.502) (9,7.798) (11,7.981) (13,8.107) (15,8.198) (17,8.267) (19,8.322) (21,8.365) (23,8.401) (25,8.432) (27,8.457) (29,8.479) (31,8.499) (33,8.516) (35,8.530) (37,8.544) (39,8.556) (41,8.566) (43,8.576) (45,8.585) (47,8.593) (49,8.601) (51,8.607) (53,8.614) (55,8.620) (57,8.625) (59,8.630) (61,8.635)};
  \foreach \p in {(1,0.000),(3,5.556),(5,6.944),(7,7.502),(9,7.798),(11,7.981),(13,8.107),(15,8.198),(17,8.267),(19,8.322),(21,8.365),(23,8.401),(25,8.432),(27,8.457),(29,8.479),(31,8.499),(33,8.516),(35,8.530),(37,8.544),(39,8.556),(41,8.566),(43,8.576),(45,8.585),(47,8.593),(49,8.601),(51,8.607),(53,8.614),(55,8.620),(57,8.625),(59,8.630),(61,8.635)} {\fill[sea] \p circle (1.1pt);}
  \node[font=\scriptsize,text=ink] at (31,-2.1) {number of voters $n$};
  \node[font=\scriptsize,text=ink,anchor=west] at (4,5.0) {$\tfrac1{18}$ at $n=3$};
\end{tikzpicture}
\caption{The probability of a Condorcet cycle under majority rule with $n$ voters, for odd $n\le61$, computed exactly from \cref{thm:fourier:kalai} and the Fourier coefficients of majority, which are known in closed form \parencite{ODonnell2014}. The value $0$ at $n=1$ is the dictator.}
\label{fig:fourier:guilbaud}
\end{figure}

The classical derivation of \cref{sec:fourier:condorcet} needed a Gaussian orthant in three dimensions; here a pair suffices, because \eqref{eq:fourier:kalai} is quadratic in $f$. More important, the formula is exact for every $n$ (\cref{fig:fourier:guilbaud}) and every rule, and it puts Condorcet, Guilbaud, and Arrow into one sentence: the probability of a cycle is $\frac14+\frac34\fourierStab_{-1/3}[f]$, which vanishes for dictators and their negations and for nothing else, and tends to $8.77\%$ for majority. It also poses a question the old setting could not attack: which rule makes cycles least likely if no voter may have much power? Kalai conjectured that majority is asymptotically optimal among rules that treat all voters alike, and proved for them the weaker bound $0.9192$ on the probability of a rational outcome (Exercise~\ref{ex:fourier:symmetric}).

\begin{theorem}[Majority is most rational; \textcite{MosselODonnellOleszkiewicz2010}]\label{thm:fourier:moo}
For every $\eps>0$ there is $\tau>0$ such that every Boolean $f$ with $f(-x)=-f(x)$ and all influences at most $\tau$, used in all three contests under impartial culture, produces a cycle with probability at least $\frac14-\frac3{2\pi}\arcsin\frac13-\eps$.
\end{theorem}

The proof starts on the spectrum: an odd function has weight only on odd levels, so $\fourierStab_{-1/3}[f]=-\fourierStab_{1/3}[f]$, and the theorem becomes a case of the ``Majority is Stablest'' bound $\fourierStab_{1/3}[f]\le\frac2\pi\arcsin\frac13+o(1)$. It then leaves the cube: an invariance principle shows that a low-influence function of many independent signs behaves like a function of Gaussians, and in Gaussian space the extremal problem was solved by Borell's isoperimetric theorem.\sidenote{Sheppard's formula reappears as the analysis of random-hyperplane rounding for maximum cut in \cref{ch:sos}; Majority is Stablest is the analytic ingredient in the proof that, under the Unique Games Conjecture, this rounding is optimal \parencite{MosselODonnellOleszkiewicz2010}.}

\section{After: the power of a single voter}\label{sec:fourier:influence}

\subsection{Influence and the Poincaré inequality}\label{sec:fourier:poincare}

For a Boolean $f$ and $i\in[n]$, write $x^{\oplus i}$ for $x$ with its $i$th coordinate flipped, so that $\fourierInf_i[f]=\Pr[f(x)\ne f(x^{\oplus i})]$, and let the \term{total influence} be $\fourierTI[f]=\sum_i\fourierInf_i[f]$.

\begin{lemma}[Influence on the spectrum]\label[lemma]{lem:fourier:inf}
Every Boolean $f$ satisfies $\fourierInf_i[f]=\sum_{S\ni i}\hat f(S)^2$ and $\fourierTI[f]=\sum_S\abs S\hat f(S)^2$.
\end{lemma}

\begin{proof}
Let $D_if(x)=\frac12\bigl(f(x^{i\to1})-f(x^{i\to-1})\bigr)$, with $x^{i\to b}$ the point $x$ with its $i$th coordinate set to $b$. On a character, $D_i$ removes $i$ from $S$ if $i\in S$ and gives $0$ otherwise, so $D_if=\sum_{S\ni i}\hat f(S)\chi_{S\setminus\{i\}}$, and Parseval's identity gives $\E[(D_if)^2]=\sum_{S\ni i}\hat f(S)^2$. On the other hand $D_if$ takes values in $\{-1,0,1\}$ and is nonzero exactly when $f(x)\ne f(x^{\oplus i})$, so $\E[(D_if)^2]=\fourierInf_i[f]$. Summing over $i$ counts each set $S$ exactly $\abs S$ times.
\end{proof}

\begin{proposition}[Poincaré inequality]\label[proposition]{prop:fourier:poincare}\index{Poincar\'e inequality}
Every Boolean $f$ satisfies $\Var f\le\fourierTI[f]$. Hence some voter has influence at least $\Var(f)/n$.
\end{proposition}

\begin{proof}
Compare the two sums term by term: $\Var f=\sum_{S\ne\emptyset}\hat f(S)^2\le\sum_S\abs S\hat f(S)^2=\fourierTI[f]$. The largest of $n$ numbers is at least their average.
\end{proof}

For a balanced rule $\Var f=1$, and the bound $1/n$ drops out in one line, without the edge-isoperimetric inequality. To go further, Parseval's identity is not enough, since it constrains only the sum of the squared coefficients; what is needed is an inequality between different norms.

\subsection{Hypercontractivity and the KKL theorem}\label{sec:fourier:kkl}

The noise operator damps level $k$ by the factor $\rho^k$, so $T_\rho f$ is smoother than $f$. Hypercontractivity\index{hypercontractivity} makes this quantitative: smoothing improves integrability.\sidenote{Hypercontractive estimates on the cube go back to \textcite{Bonami1970}. \textcite{Beckner1975} proved the sharp two-point inequality as a step toward the sharp constant in the Hausdorff--Young inequality.}

\begin{theorem}[Bonami's inequality]\label{thm:fourier:bonami}
Every real function $f$ on $\{-1,1\}^n$ satisfies $\norm{T_{1/\sqrt3}f}_4\le\norm{f}_2$.
\end{theorem}

\begin{proof}
The plan is induction on $n$, peeling off the last coordinate. For $n=0$ the function is constant and the two sides are equal. For $n\ge1$ write $f(x)=g(x')+x_nh(x')$, with $x'=(x_1,\dots,x_{n-1})$ and $g$, $h$ independent of $x_n$; then $\norm f_2^2=\norm g_2^2+\norm h_2^2$, because $\E[x_n]=0$. Let $\rho=1/\sqrt3$. Since $T_\rho$ multiplies $\chi_S$ by $\rho^{\abs S}$, we have $T_\rho f=A+\rho x_nB$ with $A=T_\rho g$ and $B=T_\rho h$, both computed on $n-1$ coordinates. Expand the fourth power and average over $x_n$ first: the odd powers of $x_n$ vanish, $6\rho^2=2$, and $\rho^4=\frac19\le1$. Bounding the middle term by the Cauchy--Schwarz inequality,
\[
  \E[(T_\rho f)^4]=\E[A^4]+2\E[A^2B^2]+\tfrac19\E[B^4]\le\bigl(\norm A_4^2+\norm B_4^2\bigr)^2 .
\]
The induction hypothesis bounds the right side by $\bigl(\norm g_2^2+\norm h_2^2\bigr)^2=\norm f_2^4$.
\end{proof}

\begin{corollary}[The dual form]\label[corollary]{cor:fourier:dualbonami}
Every real function $f$ on $\{-1,1\}^n$ satisfies $\norm{T_{1/\sqrt3}f}_2\le\norm{f}_{4/3}$.
\end{corollary}

\begin{proof}
Write $T$ for $T_{1/\sqrt3}$, which is self-adjoint by \eqref{eq:fourier:noisecorr}. By Hölder's inequality with exponents $\frac43$ and $4$, and then \cref{thm:fourier:bonami} applied to $Tf$,
\[
  \norm{Tf}_2^2=\ip{f,T(Tf)}\le\norm f_{4/3}\norm{T(Tf)}_4\le\norm f_{4/3}\norm{Tf}_2 .
\]
Divide by $\norm{Tf}_2$; if it is $0$, there is nothing to prove.
\end{proof}

\begin{theorem}[\textcite{KahnKalaiLinial1988}]\label{thm:fourier:kkl}\index{Kahn--Kalai--Linial theorem}
For every $n\ge2$, every Boolean function $f$ on $\{-1,1\}^n$ has a voter whose influence is at least $\Var(f)\ln n/(9n)$.
\end{theorem}

\begin{proof}
The plan is to play two facts against each other: when every influence is small, hypercontractivity makes the low levels of the spectrum light; when the total influence is small, the spectrum must sit on low levels. Write $V=\Var f$, $\fourierTI=\fourierTI[f]$, and $M=\max_i\fourierInf_i[f]$. If $V=0$, then there is nothing to prove, so assume $V>0$; then $\fourierTI\ge V>0$ by \cref{prop:fourier:poincare}.

\emph{Small influences make the low levels light.} The derivative $D_if$ of \cref{lem:fourier:inf} takes values in $\{-1,0,1\}$ and is nonzero with probability $\fourierInf_i[f]$, so $\norm{D_if}_{4/3}^2=\fourierInf_i[f]^{3/2}$. Apply \cref{cor:fourier:dualbonami} to $D_if=\sum_{S\ni i}\hat f(S)\chi_{S\setminus\{i\}}$ and sum over $i$:
\begin{equation}\label{eq:fourier:kkl1}
  \sum_S\abs S\,3^{1-\abs S}\hat f(S)^2=\sum_i\norm{T_{1/\sqrt3}D_if}_2^2\le\sum_i\fourierInf_i[f]^{3/2}\le\sqrt M\,\fourierTI .
\end{equation}

\emph{Small total influence makes the low levels heavy.} Let $k=\ceil{2\fourierTI/V}$. The weight above level $k$ is at most $\fourierTI/(k+1)<V/2$, so the levels $1,\dots,k$ carry more than $V/2$. On these levels $\abs S\,3^{1-\abs S}\ge3^{1-k}$, and \eqref{eq:fourier:kkl1} gives
\begin{equation}\label{eq:fourier:kkl2}
  V<2\cdot3^{k-1}\sqrt M\,\fourierTI .
\end{equation}

\emph{The two cases.} If $\fourierTI\ge V\ln n/9$, then the largest influence is at least the average, $\fourierTI/n\ge V\ln n/(9n)$. Otherwise $k-1<2\fourierTI/V<\frac29\ln n$, so $3^{k-1}<n^{(2\ln3)/9}<n^{1/4}$, and \eqref{eq:fourier:kkl2} gives $\sqrt M>V/(2n^{1/4}\fourierTI)>9/(2n^{1/4}\ln n)$. Squaring, $M>81/(4\sqrt n\ln^2n)$. This is at least $\ln n/(9n)$, which is at least $V\ln n/(9n)$, because $729\sqrt n\ge4\ln^3n$ for all $n\ge2$ (the ratio $\ln^3n/\sqrt n$ never exceeds $11$).
\end{proof}

The tribes function shows that the theorem is sharp up to the constant (Exercise~\ref{ex:fourier:tribes}). The proof lives in the Fourier world twice over: smoothing, which \cref{thm:fourier:bonami} controls, acts level by level only because the noise operator is diagonal, and the argument splits $f$ by level, which has no meaning in terms of the values of $f$.

\section{Costs and limits}\label{sec:fourier:costs}

\paragraph{Where did the difficulty go?}
Into two places of very different size. The first is the dictionary, which is cheap (two lines for orthonormality, one for the eigenvalues of the noise operator) because the problems had already paid for it: the voters are independent and each ballot is as likely as its reverse, and the queries of the test are independent and uniform. Product structure and a group action make each experiment invariant, and invariance makes the operators diagonal. The move pays once, when one recognizes the probability in question as a quadratic or cubic form with an invariant kernel; afterward each application is a computation of eigenvalues. The second place is analytic. Parseval's identity says only that the squared coefficients sum to $1$, and every result that needs more needs an inequality between different norms. Hypercontractivity is the basic one, paid for once by the induction in \cref{thm:fourier:bonami} and reused in KKL, in the invariance principle behind \cref{thm:fourier:moo}, and, with its reverse form, in the quantitative Arrow theorems of \textcite{Mossel2012} and \textcite{Keller2012}. Sharp constants cost more still: the optimality of majority needs the invariance principle.

\paragraph{What did not dissolve.}
Everything that breaks the symmetry. If voters are correlated, as members of a party are, then the noise operator stops being diagonal and Kalai's formula acquires cross terms between different sets $S$. On a slice of the cube, the profiles with a fixed number of votes for each candidate, the characters are not orthogonal; on a non-abelian group they give way to irreducible representations, and invariant operators become only block diagonal. Arrow's theorem with three different rules is not reached by the spectral argument alone (\cref{rem:fourier:threerules}), and its quantitative version needed much heavier tools. Exact constants often escape: the bound $\Pr[\text{reject}]\ge\dist(f,\fourierLin)$ is sharp only near distance $\frac12$, for small distances the truth is $3d-6d^2$ (Exercise~\ref{ex:fourier:corrupt}), and the minimum rejection rate $\frac{45}{128}$ at distances of at least $\frac14$ came, in the same paper of \textcite{BellareEtAl1996}, from a combinatorial isoperimetric argument. Fourier coefficients control configurations of three points, such as the triple $(x,y,x\odot y)$ or a three-term progression, but not of four, and \textcite{Gowers2001} needed a higher-order theory for longer progressions. Finally, the spectrum is a structure, not an algorithm: a large coefficient need not be easy to find from random examples, and for parities observed with random noise no polynomial-time learning algorithm is known.

\paragraph{Other seas.}
Gaussian space, where Hermite polynomials diagonalize the Ornstein--Uhlenbeck semigroup, is the limit of the cube under the central limit theorem; Sheppard's formula lives there, and the invariance principle carries extremal results back. Representation theory extends the move to non-abelian groups. Hypercontractivity is equivalent to a logarithmic Sobolev inequality, a link to the transport methods of \cref{ch:transport}.

\section{The move, distilled}\label{sec:fourier:distilled}

\begin{move}{Diagonalize the symmetry}
\movepart{In one sentence} When the random experiment behind a problem is invariant under a commutative group, expand the unknown function in the characters of the group: every invariant operator becomes diagonal, and the probability you want becomes a weighted sum over the spectrum.
\movepart{Triggers} The quantity of interest is the expectation of a polynomial in the values of one function at a few correlated random points whose law is invariant under a common translation. The operators in sight are averages, noise, or convolutions; the conjectured extremal functions are characters, dictators, or majority; the question is how close to linear, or to a dictator, a function with a given statistic must be.
\movepart{Recipe} Identify the group and check that the experiment is invariant. Write down the characters and prove them orthonormal. Compute the eigenvalue of each operator on each character; by \cref{lem:fourier:diag} there is nothing else to compute. Rewrite the target as a sum over the spectrum, use Parseval's identity as the constraint that Booleanness imposes, and optimize level by level. Translate the extremal spectrum back into structure: weight on level $1$ means a dictator. When Parseval's identity is not enough, add an inequality between norms.
\movepart{Anti-pattern} Expanding in characters when the experiment is not invariant, as with correlated voters or a non-product measure: the operators are not diagonal, every computation produces cross terms, and the expansion is mere notation. So is a combinatorial proof rewritten in Fourier coefficients that never uses an eigenvalue.
\end{move}

\dissolutiontest{Linearity test: three lines replace self-correction, with constant $1$ for $2$. Arrow for one rule: one identity and one inequality replace hand-built profiles.}{The test analysis holds at every rejection rate; Kalai's formula, for every rule and every $n$.}{Condorcet's paradox, Guilbaud's $8.77\%$, and Arrow's theorem are three readings of \eqref{eq:fourier:kalai}; tests and elections are diagonal forms.}{The KKL theorem; quantitative Arrow theorems; majority is most rational; optimal inapproximability.}

The move is one of several ways to use a symmetry that a problem has for free. \Cref{ch:conformal} uses exchangeability through invariance alone, with no basis, and the symmetrization of \cref{ch:vc} averages over a swap symmetry; here the symmetry chooses the coordinates, and the price is that it must be exact. The move also has a rival: for sets in $\F_3^n$ without three-term progressions, the Fourier density increment of Exercise~\ref{ex:fourier:meshulam} gives $O(3^n/n)$, and the polynomial method of \cref{ch:dimension} an exponentially smaller bound. Eigenvalues of Cayley graphs (Exercise~\ref{ex:fourier:cayley}) are the expander side of \cref{ch:pseudorandom}.

\begin{exercises}
\item \consolidation (Three and five voters.) (a) List the $12$ profiles of three voters that produce a cycle under majority. (b) Show that $\widehat{\fourierMaj_5}(S)$ is $\frac38$ for $\abs S=1$, $-\frac18$ for $\abs S=3$, $\frac38$ for $\abs S=5$, and $0$ otherwise, and check Parseval's identity. (c) Deduce from \cref{thm:fourier:kalai} that five voters produce a cycle with probability $\frac5{72}$.\label{ex:fourier:maj5}

\item \consolidation (A corrupted character.) Let $n\ge2$, let $S\subseteq[n]$, and let $f$ agree with $\chi_S$ except at one point $p\ne\fourierOnes$. Show that the linearity test rejects $f$ with probability exactly $3d-6d^2$, where $d=2^{-n}=\dist(f,\fourierLin)$, so that the bound of \cref{thm:fourier:blr} is off by a factor close to $3$ for small distances. What is the rejection probability if $p=\fourierOnes$?\label{ex:fourier:corrupt}

\item \consolidation (Poincaré and edge isoperimetry.) (a) Show that a Boolean $f$ satisfies $\Var f=\fourierTI[f]$ if and only if it is constant, a dictator, or the negation of a dictator. (b) Let $A\subseteq\{-1,1\}^n$ have $\alpha2^n$ points. Apply \cref{prop:fourier:poincare} to the function that is $-1$ on $A$ and $1$ off $A$ to show that at least $2\alpha(1-\alpha)2^n$ edges of the cube join $A$ to its complement. Which sets attain equality when $\alpha=\frac12$?\label{ex:fourier:poincare}

\item \consolidation\hard (Rules that treat voters alike.) Call $f$ transitive-symmetric if for all voters $i$ and $j$ some permutation of the coordinates maps $i$ to $j$ and leaves $f$ unchanged. (a) Show that such $f$ satisfy $\fourierW1[f]\le\E\bigl[\abs{x_1+\dots+x_n}\bigr]^2/n$, with equality for majority, and that the right side tends to $2/\pi$. (b) Deduce from \cref{thm:fourier:arrowq} that such a rule, used in all three contests, gives a rational outcome with probability at most $\frac79+\frac4{9\pi}+o(1)\approx0.9192$, Kalai's bound; compare with majority's $0.9123$. \hint{in (a) the coefficients $\hat f(\{i\})$ are equal, and their sum is $\E[f(x)(x_1+\dots+x_n)]$.}\label{ex:fourier:symmetric}

\item \consolidation\hard (Arrow's theorem with three rules.) Let $(f,g,h)$ be unanimous and rational. (a) Show that $f(x)=-h(-x)$ for every $x$. (b) Relabel the candidates cyclically to deduce that $f=g=h$ and that $f$ is odd. (c) Conclude \cref{thm:fourier:arrow} from \cref{thm:fourier:arrowq}. \hint{use the profiles in which every voter ranks $c$ immediately below $b$, and those in which every voter ranks $c$ immediately above $b$; in both, the contest between $a$ and $c$ copies the contest between $a$ and $b$.}\label{ex:fourier:neutral}

\item \consolidation (Tribes.) Split $n=ws$ voters into $s$ blocks of $w$, and let $f(x)=1$ if some block votes unanimously $+1$ and $f(x)=-1$ otherwise. (a) Show that every voter has influence $2^{1-w}(1-2^{-w})^{s-1}$. (b) With $s=\ceil{2^w\ln2}$, show that $\Pr[f=1]\to\frac12$ and that every influence is $(1+o(1))\ln n/n$ as $w\to\infty$. Conclude that \cref{thm:fourier:kkl} is sharp up to the constant factor.\label{ex:fourier:tribes}

\item \transfer (Cayley graphs.) Let $G$ be a finite abelian group and $A=-A$ a subset; the Cayley graph joins $u$ to $u+a$ for $a\in A$. (a) Show that the adjacency operator $(Kf)(u)=\sum_{a\in A}f(u+a)$ commutes with translations and has each character $\chi$ of $G$ (a homomorphism to the unit circle) as an eigenvector, with eigenvalue $\sum_{a\in A}\chi(a)$. (b) For the cube, $G=\F_2^n$ with $A$ the $n$ unit vectors, show that the eigenvalues are $n-2k$ with multiplicity $\binom nk$. (c) For the cycle, $G=\Z/m\Z$ with $m\ge3$ and $A=\{\pm1\}$, show that they are $2\cos(2\pi j/m)$ for $j=0,\dots,m-1$.\label{ex:fourier:cayley}

\item \transfer\harder (Meshulam's theorem.) Let $A\subseteq\F_3^n$ have density $\alpha$ and no three distinct points $x$, $y$, $z$ with $x+y+z=0$, that is, no three-term progression. Let $\hat A(r)=\E_x[\ind_A(x)\omega^{-r\cdot x}]$ for $r\in\F_3^n$, with $\omega=e^{2\pi i/3}$. (a) Show that $\sum_r\hat A(r)^3=\alpha3^{-n}$, and deduce that $\abs{\hat A(r)}\ge\alpha^2/2$ for some $r\ne0$ when $3^n\ge2/\alpha^2$. (b) Show that $A$ then has density at least $\alpha+\frac38\alpha^2$ on one of the three hyperplanes $\{x\colon r\cdot x=j\}$, and that its part there, translated into $\{x\colon r\cdot x=0\}\cong\F_3^{n-1}$, again contains no progression. (c) Iterate to show that $\abs A=O(3^n/n)$, the bound of \textcite{Meshulam1995}. \hint{in (a), the only solutions with $x,y,z\in A$ are the $\abs A$ triples $x=y=z$; in (b), write $\hat A(r)=\frac13\sum_j(a_j-\alpha)\omega^{-j}$, with $a_j$ the density of $A$ on the $j$th hyperplane.}\label{ex:fourier:meshulam}

\item \transfer\hard (Learning from the spectrum.) Let $f$ be Boolean with $\sum_{\abs S>d}\hat f(S)^2\le\eps$, and let $N=\sum_{k\le d}\binom nk$. (a) Show that $\Pr[f(x)\ne\sgn g(x)]\le\E[(f(x)-g(x))^2]$ for every real $g$, with $\sgn0=1$. (b) Suppose that $\abs{c_S-\hat f(S)}\le\eta$ for $\abs S\le d$, and let $g=\sum_{\abs S\le d}c_S\chi_S$. Show that $\sgn g$ errs on a fraction at most $\eps+N\eta^2$ of the inputs. (c) Show that $O\bigl((N/\eps)\log(N/\beta)\bigr)$ examples $(x,f(x))$ with uniform $x$ make this fraction at most $2\eps$ with probability $1-\beta$. By \textcite{LinialMansourNisan1993}, circuits of constant depth and size $s$ satisfy the hypothesis with $d$ polylogarithmic in $s/\eps$; deduce a learning algorithm for them in time $n^{\operatorname{polylog}(s/\eps)}\log(1/\beta)$. \hint{use Hoeffding's inequality and a union bound, with $\eta^2=\eps/N$.}\label{ex:fourier:lmn}

\item \research\harder (A basis for neural networks.) For $f$ on $\{-1,1\}^d$ with Fourier support $\{S_1,\dots,S_m\}$, \textcite{AbbeBoixAdseraMisiakiewicz2023} define the leap as the minimum over orderings of the support of the largest number of new coordinates that one set adds to the union of the earlier ones. For uniform Boolean and for Gaussian inputs, they conjecture that stochastic gradient descent on fully connected networks learns a function of boundedly many coordinates in time $d^{\max(\mathrm{leap},2)}$ up to logarithmic factors, and they prove a version for two-layer networks in some Gaussian cases, with modified training. (a) Compute the leap of $\chi_{[k]}$, of $\fourierMaj_3$, and of $x_1+x_1x_2+x_1x_2x_3$. (b) Open: prove the conjecture for Boolean inputs, or find the analogue of the leap for non-uniform inputs, where characters no longer diagonalize the relevant operators. \hint{for (b), start from the discussion in their paper.}\label{ex:fourier:leap}
\end{exercises}

\notesandrefs
Our treatment of the linearity test, of Kalai's formula, and of the KKL theorem follows the route of \textcite{ODonnell2014}, with the constants worked out; the book carries each topic much further.

\Textcite{Geanakoplos2005} gives three short combinatorial proofs of Arrow's theorem. The paper of \textcite{Guilbaud1952} contains, besides the limit $0.0877$, a dictatorship theorem for the aggregation of judgments proved implicitly with an ultrafilter, the device of our aside. The quantitative theorems that followed \textcite{Kalai2002} are surveyed in the introduction of \textcite{Keller2012}, and the history of the testing constants in \textcite{BellareEtAl1996}. For longer progressions, \textcite{Gowers2001} explains why characters no longer suffice.

\printchapterbib
\end{refsection}
